\chapter{Lattice Fermions And Chiral Symmetry}
\label{ch:constructive-lattice-fermions-chiral-symmetry}

\ConventionsEuclidean

Fermions on a lattice require more structure than scalar fields.  The
variables in the path integral are Grassmann-valued algebraic variables, the
operator entering the action is a finite matrix at finite lattice volume, and
reflection positivity depends on how the lattice is cut.  This chapter
records the core fermion constructions and the chiral-symmetry problem.

\section{Finite Lattice Grassmann Integral}

Let \(\Lambda\subset a\mathbb Z^D\) be finite.  A lattice fermion field is a
collection of generators
\[
  \psi_\alpha(x),\qquad \bar\psi_\alpha(x),
  \qquad x\in\Lambda ,
\]
of a finite Grassmann algebra.  For a matrix \(D_{xy}^{\alpha\beta}\), the
Gaussian integral is the Berezin functional defined by
\[
  \int \prod_{x,\alpha}d\bar\psi_\alpha(x)d\psi_\alpha(x)\,
  \exp\left(
    -\sum_{x,y}\bar\psi_\alpha(x)D_{xy}^{\alpha\beta}\psi_\beta(y)
  \right)
  =
  \det D .
\]
Correlation functions are algebraic coefficients of this finite Berezin
integral.  No Borel measure on a set of fermion configurations is involved.

\section{Naive Dirac Operator And Doublers}

The naive lattice Dirac operator is
\[
  D_{\mathrm{naive}}(p)
  =
  \frac{i}{a}\sum_{\mu=1}^D\gamma_\mu\sin(ap_\mu).
\]
It vanishes not only at \(p=0\), but also at every corner
\[
  p_\mu\in\{0,\pi/a\}.
\]
Thus it describes \(2^D\) light fermion species in the continuum expansion.
The doubling is a statement about the zeros and chiralities of the lattice
symbol.  Removing it while keeping locality, translation invariance,
reflection positivity, and exact chiral symmetry is the central obstruction.

\section{Wilson Fermions}

Wilson's operator adds a lattice Laplacian:
\[
  D_W(p)
  =
  \frac{i}{a}\sum_\mu\gamma_\mu\sin(ap_\mu)
  +
  m
  +
  \frac{r}{a}\sum_\mu(1-\cos(ap_\mu)).
\]
The Wilson term gives masses of order \(a^{-1}\) to the doubler modes.  It
also breaks the naive chiral symmetry at nonzero lattice spacing.  The
continuum chiral symmetry is recovered only after tuning the mass and taking
the continuum limit in a regime where the target theory exists.

With gauge fields \(U_\mu(x)\), the Wilson operator is obtained by replacing
finite differences with covariant finite differences.  Gauge invariance is
exact at finite lattice spacing, while chiral symmetry is not.

\section{Reflection Positivity}

Let the time reflection be \(\theta:x_0\mapsto -x_0+a\), so that the
reflection plane lies between two time slices.  For Wilson fermions the
action can be decomposed as
\[
  S=S_+ + \theta S_+ + S_{\mathrm{mid}},
\]
where \(S_{\mathrm{mid}}\) contains only hopping terms crossing the mid-link
reflection plane.  Reflection positivity follows when the mid-link term can
be written as a finite sum
\[
  S_{\mathrm{mid}}=-\sum_i \theta(F_i)F_i
\]
with the correct Grassmann reflection signs and with positive gauge-link
weights.  The precise reflection operation includes complex conjugation,
time reflection of links, and the spin matrix implementing Euclidean time
reflection.

\section{Ginsparg--Wilson Relation}

A lattice Dirac operator \(D\) satisfies the Ginsparg--Wilson relation when
\[
  \gamma_5D+D\gamma_5=aD\gamma_5D .
\]
Then the action \(\bar\psi D\psi\) is invariant under the lattice chiral
transformation
\[
  \delta\psi=\gamma_5(1-aD)\psi,\qquad
  \delta\bar\psi=\bar\psi\gamma_5 .
\]
The fermion measure changes by a finite Berezinian whose logarithm gives the
lattice index
\[
  \operatorname{index}(D)
  =
  \operatorname{Tr}\,\gamma_5\left(1-\frac{a}{2}D\right).
\]
This is the lattice version of the axial anomaly in a gauge background.

\section{Overlap Operator}

Let
\[
  H_W=\gamma_5(D_W-m_0)
\]
with \(m_0\) chosen in the physical range.  The overlap Dirac operator is
\[
  D_{\mathrm{ov}}
  =
  \frac{1}{a}\left(1+\gamma_5\operatorname{sign}(H_W)\right).
\]
It satisfies the Ginsparg--Wilson relation when \(H_W\) has no zero modes.
Locality requires gauge fields to lie in an admissible region or to be
controlled by estimates on the spectrum of \(H_W\).  Thus exact lattice
chiral symmetry is obtained at the price of a non-ultralocal operator whose
locality is conditional.

\section{Continuum Ledger}

A lattice fermion construction must state:
\[
\begin{gathered}
  \text{finite Grassmann algebra and Berezin integral},\\
  \text{Dirac operator, gauge covariance, and boundary conditions},\\
  \text{reflection positivity or a replacement Hilbert-space construction},\\
  \text{chiral symmetry, anomaly, and index convention},\\
  \text{continuum scaling window and operator renormalization}.
\end{gathered}
\]
For chiral gauge theories, the determinant or Pfaffian is a section of a
determinant line over gauge fields.  A nonperturbative definition requires a
gauge-invariant choice of phase compatible with locality and anomaly
cancellation.

\begin{openproblem}[Chiral gauge lattice]
\label{op:chiral-gauge-lattice}
Construct four-dimensional interacting chiral gauge theories on the lattice
with locality, reflection positivity or an equivalent reconstruction route,
exact gauge invariance, anomaly cancellation, and a continuum limit with the
desired local operator algebra.
\end{openproblem}
